% ============================================================
%  part07.tex  |  原文 PDF p.48–55（页脚 p.47–54）
%  内容：第 6 章 BRST–BV 形式（BRST-BV Formalism）
%        6.1 商与鬼场（Quotient and Ghost）
%        6.2 多向量场（Polyvector Fields）
%        6.3 Batalin–Vilkovisky（BV）代数
%        6.4 经典 BV 主方程（Classical BV Master Equation，起于 PDF p54）
%  说明：6.5 量子主方程起于 PDF p56，属 part08；本片段止于 p55。
% ============================================================

% >>> p048

\section{BRST--BV 形式（BRST-BV Formalism）}

\begin{keypoints}
  \item \textbf{核心物理图像}：规范理论的轨道空间 $V/G$ 往往奇异，直接在 $V/G$ 上
        积分一般不可行。两条路线是：Faddeev--Popov 型（在 $V$ 上取规范切片 $Y$，
        并以 $\delta(Y)\xi_Y$ 配重使积分在 $Y$ 的连续形变下不变）与 BRST--BV 型
        （用导出商 $V\times\Lie{g}[-1]$ 代替 $V/G$，鬼场正是导出方向的坐标）。
  \item \textbf{数学工具}：Chevalley--Eilenberg 复形
        $C^{\bullet}(\Lie{g},\mathcal{O}(V))$（BRST 微分 $d_{\mathrm{CE}}$）、
        多向量场 $\mathrm{PV}(M)$（即 $\Gamma(M,\wedge^{\bullet}TM)$）、由体积形式诱导的
        BV 算子 $\Delta$（广义散度 / 奇 Laplacian），以及 BV 括号
        （Schouten--Nijenhuis 括号的奇类比）。
  \item \textbf{主要结论}：在 BV 代数 $(A,\Delta)$ 中，
        $\{a,b\}:=\Delta(ab)-(\Delta a)b-(-1)^{|a|}a\Delta b$ 度量 $\Delta$ 偏离
        导子的程度；若 $S_0$ 满足经典主方程 $\{S_0,S_0\}=0$，则
        $\delta=\{S_0,-\}$ 满足 $\delta^2=0$，给出经典 BRST 微分；任何 $G$-不变
        经典作用量都可自然延拓为 CME 的解（BV 形式的普适性）。
  \item \textbf{译者评注}：本章是全书技术核心之一，难点集中在\emph{奇偶符号}、
        $\Delta$ 的归一化与 Schouten 括号的符号约定；6.4 节只给出经典主方程，
        量子主方程（QME）与相关上同调的解释留待第 6.5 节（part08）。
\end{keypoints}

我们讨论经典与量子规范理论的 BRST--BV 形式。本节我们聚焦于有限维玩具模型，
以说明主要结构。下一节我们讨论如何纳入重整化。

\subsection{商与鬼场（Quotient and Ghost）}

考虑如下有限维规范理论玩具模型
\[
  f:\ V\longrightarrow\C,
\]
其中连通 Lie 群 $G$ 作用在 $V$ 上，且 $f$ 是 $G$-不变的。等价地，$f$ 是 $V/G$
上的函数，而我们关心轨道空间上的积分
\[
  \int_{V/G}\ee^{\frac{1}{\hbar}f}(-).
\]
通常有两种处理方式。

\textcircled{1}\ 寻找一个子空间 $Y\subset V$，使得 $Y$ 中的点代表 $G$ 的轨道，
并构造一个与 $Y$ 相联系的“形式” $\xi_Y$，使得
\[
  \int_V \ee^{\ii f/\hbar}\delta(Y)\xi_Y(-)
\]
在 $Y$ 的连续形变下保持不变。这里的 $Y$ 称为\emph{规范固定条件}。著名的
Faddeev--Popov 方法即属此类。
\begin{figure}[htbp]
  \centering
  \includegraphics[width=0.52\linewidth]{fig_p048_1.png}
  \caption*{规范固定切片：$Y\subset V$（原文 Figure 2，原文 p.47 / PDF p.48）}
\end{figure}

% [Tutor Note: 两种方案的几何区别要抓住：Faddeev–Popov 仍在原空间 $V$ 上积分，
% 只是插入 $\delta(Y)\xi_Y$ 把积分“压”到切片 $Y$ 上，并证明结果不随 $Y$ 的连续
% 形变而变（这相当于在 $V$ 上构造了一个闭形式）；BRST–BV 则改用导出商
% $V\times\Lie{g}[-1]$，从而在轨道空间 $V/G$ 本身有奇点时仍然有效。原文
% $\delta(Y)\xi_Y$ 里的 $\xi_Y$ 是与 $Y$ 的选取相关的伴随“形式”，作者未写出其
% 具体表达式，此处不宜擅自补全。]

当商 $V/G$ 的奇异性不强时，上述构造效果很好。一般地，若轨道空间 $V/G$ 是
“坏”的，就会出现新问题。这是几何中的典型问题，而教训是我们应当取导出对象
来消解奇异性。这引出了另一种方法。

\textcircled{2}\ BRST--BV 方法。令 $\mathcal{O}(V)$ 表示 $V$ 上的函数。则 $V/G$
上的函数
\[
  \mathcal{O}(V/G)=\mathcal{O}(V)^G
\]
可以自然地等同于 $V$ 上的 $G$-不变函数。假设 $G$ 是连通 Lie 群，Lie 代数为
$\Lie{g}$。每个 $u\in\Lie{g}$ 生成 $V$ 上的一个向量场 $X_u$，而 $G$-不变函数
可以无穷小地表示为
\[
  \mathcal{O}(V)^G=\{\phi\in\mathcal{O}(V)\mid X_u(\phi)=0,\ \forall u\in\Lie{g}\}.
\]

% >>> p049

设 $\{e_1,\dots,e_m\}$ 为 $\Lie{g}$ 的一组基，
$\{c^1,\dots,c^m\}$ 为 $\Lie{g}^{\vee}$ 的对偶基。记 $X_\alpha:=X_{e_\alpha}$。
回顾 Chevalley--Eilenberg 微分（BRST）
\[
  \begin{aligned}
    d_{\mathrm{CE}}:\ \mathcal{O}(V)&\longrightarrow
      \Lie{g}^{\vee}\otimes\mathcal{O}(V),\\
    \phi&\longmapsto\sum_\alpha c^\alpha\otimes X_\alpha(\phi).
  \end{aligned}
\]
我们有如下自然的正合列
\[
  0\longrightarrow\mathcal{O}(V)^G\longrightarrow\mathcal{O}(V)
  \longrightarrow\Lie{g}^{\vee}\otimes\mathcal{O}(V).
\]
现在清楚了：完整的 Chevalley--Eilenberg 复形
\[
  C^{\bullet}(\Lie{g},\mathcal{O}(V))
  =\left(\wedge^{\bullet}\Lie{g}^{\vee}\otimes\mathcal{O}(V),
  d_{\mathrm{CE}}\right)
\]
是 $\mathcal{O}(V)^G$ 的一个自然的导出替代（derived replacement）。
$\Lie{g}^{\vee}$ 中的反交换变量 $c^i$ 称为\emph{鬼场}，如此前所见。由于
\[
  \wedge^{\bullet}\Lie{g}^{\vee}
  =\Sym^{\bullet}\left(\Lie{g}^{\vee}[-1]\right)
  =\Sym^{\bullet}\left((\Lie{g}[1])^{\vee}\right),
\]
我们可以把 $\wedge^{\bullet}\Lie{g}^{\vee}$ 看作（奇）分次空间 $\Lie{g}[-1]$
上的自由多项式环。于是
\[
  C^{\bullet}(\Lie{g},\mathcal{O}(V))=
  V\times\Lie{g}[-1]\ \text{上的函数}.
\]
这是一个 dg 空间，因为 $C^{\bullet}(\Lie{g},\mathcal{O}(V))$ 是以
Chevalley--Eilenberg 微分为微分的 dg 代数
\[
  d_{\mathrm{CE}}
  =\frac{1}{2}f^{\gamma}_{\ \alpha\beta}c^{\alpha}c^{\beta}
  \frac{\partial}{\partial c^{\gamma}}+c^{\alpha}X_{\alpha}.
\]
我们应把 $V\times\Lie{g}[-1]$ 看作 $V$ 的导出 $G$-商。我们的下一个目标是在
$V\times\Lie{g}[-1]$ 上建立积分理论
\[
  \int_{V\times\Lie{g}[1]}\ee^{f/\hbar},
\]
使其能够自然推广到 QFT 的 $\infty$ 维情形。

% [Tutor Note: 上面这个正合列只在 $H^0$ 处给出不变函数：连通 $G$ 时
% $H^0(\Lie{g}^{\vee}\otimes\mathcal{O}(V))=\mathcal{O}(V)^G$，但右端映射一般不是
% 满射，$H^1$ 度量的是“不能由单个无穷小对称生成的规范冗余”。原文称其为
% “natural exact sequence” 是简写，严格地说应读作上链复形的正合列的一部分。]

% [Tutor Note: 这里两次出现 $\Lie{g}[-1]$，最后一行积分却写成
% $\int_{V\times\Lie{g}[1]}$，前后记号不一致，疑为原文笔误（应为
% $V\times\Lie{g}[-1]$：鬼场 $c^\alpha$ 的次数为 $+1$，作为坐标它属于
% $\Lie{g}[-1]$）。按体例照录原文，请核对原书 p.48 / PDF p.49。]

% [Tutor Note: $d_{\mathrm{CE}}=\frac{1}{2}f^{\gamma}_{\ \alpha\beta}
% c^{\alpha}c^{\beta}\partial_{c^{\gamma}}+c^{\alpha}X_{\alpha}$ 中的 $\frac{1}{2}$
% 依赖结构常数对 $(\alpha,\beta)$ 反对称 $f^{\gamma}_{\ \alpha\beta}
% =-f^{\gamma}_{\ \beta\alpha}$。若某文献的结构常数用不同的对称化约定，这个
% $\frac{1}{2}$ 与整体符号都要相应调整；这是核对 BRST 公式时最常见的差异来源。]

\subsection{多向量场（Polyvector Fields）}

现在我们讨论关于积分的另一种观点。设 $M$ 是 $n$ 维定向流形（或简单地取
$\R^n$）。设 $\Omega$ 是 $M$ 上的体积形式。它定义了
\[
  \begin{aligned}
    \int:\ C^{\infty}(M)&\longrightarrow\R,\\
    f&\longmapsto\int_M\Omega f.
  \end{aligned}
\]
几何上，体积形式 $\Omega$ 诱导出散度算子
\[
  \operatorname{Div}:\ \operatorname{Vect}(M)\longrightarrow C^{\infty}(M),
\]
其定义为
\[
  \int\Omega\operatorname{Div}(X)\rho
  =-\int\Omega X(\rho),\qquad \forall\rho\in C_c^{\infty}(M).
\]
具体地，
\[
  \operatorname{Div}(X)=\frac{\mathcal{L}_X\Omega}{\Omega}
  =\frac{\dd\iota_X\Omega}{\Omega}.
\]

% [Tutor Note: 这里散度的符号约定是全局性的：$\int\Omega\,\mathrm{Div}(X)\rho
% =-\int\Omega\,X(\rho)$ 中的负号来自分部积分（无边界项，因 $\rho$ 紧支集），
% 它决定了 $\mathrm{Div}=\Delta|_{\mathrm{PV}^1}$ 的符号，进而决定整个 BV 算子
% $\Delta$ 的符号。若把 $\Delta$ 与 $\mathrm{Div}$ 差一个负号，后面的 CME/QME
% 都会跟着变号；核对时务必先固定这一约定。]

% >>> p050

\begin{wdefn}{6.1}
我们定义\emph{多向量场}
\[
  \mathrm{PV}(M)=\bigoplus_k \mathrm{PV}^k(M),
  \qquad\text{其中}\qquad
  \mathrm{PV}^k(M):=\Gamma(M,\wedge^k TM),\quad k\ge 0.
\]
对 $\alpha\in\mathrm{PV}^k(M)$，记 $|\alpha|=k$。
\end{wdefn}

注意微分形式是
\[
  \Omega^k(M)=\Gamma(M,\wedge^k T^{\vee}_M).
\]
因此多向量场可以看作微分形式的对偶概念。$\mathrm{PV}(M)$ 带有类似的外积
代数结构
\[
  \mathrm{PV}^k(M)\otimes\mathrm{PV}^l(M)
  \xrightarrow{\ \wedge\ }\mathrm{PV}^{k+l}(M),
\]
它是分次交换的
\[
  \alpha\wedge\beta=(-1)^{|\alpha||\beta|}\beta\wedge\alpha.
\]
当 $k=1$ 时，$\mathrm{PV}^1(M)=\operatorname{Vect}(M)$ 正是向量场。

体积形式 $\Omega$ 让我们可以用下式把多向量场与微分形式等同起来
\[
  \begin{aligned}
    \Gamma:\ \mathrm{PV}^k(M)&\xrightarrow{\ \Omega\ }\Omega^{n-k}(M),\\
    \mu&\longmapsto\mu\lrcorner\,\Omega.
  \end{aligned}
\]
这里对 $\mu=X_1\wedge X_2\wedge\cdots\wedge X_k$，$X_i\in\operatorname{Vect}(M)$，
\[
  \mu\lrcorner\,\Omega=\iota_{X_1}\iota_{X_2}\cdots\iota_{X_k}\Omega.
\]
由于 $\Omega$ 是体积形式，$\Gamma$ 给出线性同构
$\mathrm{PV}^k\overset{\simeq}{\longrightarrow}\Omega^{n-k}$。

我们可以用 $\Gamma$ 把 de Rham 复形 $(\Omega^{\bullet}(M),\dd)$ 变换到
$\mathrm{PV}(M)$ 上：
\begin{figure}[htbp]
  \centering
  \includegraphics[width=0.61\linewidth]{fig_p050_1.png}
  \caption*{de Rham 复形与多向量场复形之间的同构（原文 PDF p.50 示意图，原文无编号）}
\end{figure}
验证：$\Delta:\mathrm{PV}^1(M)=\operatorname{Vect}(M)\to\mathrm{PV}^0(M)
=C^{\infty}(M)$ 正是与 $\Omega$ 相联系的散度算子。一般地，
\[
  \Delta:\ \mathrm{PV}^k(M)\longrightarrow\mathrm{PV}^{k-1}(M)
\]
可以看作多向量场上的广义散度算子。注意 $\dd^2=0$ 蕴含
\[
  \Delta^2=0.
\]
在物理学中，$\Delta$ 称为 \emph{Batalin--Vilkovisky（BV）算子}。

为看出 $\Delta$ 的模样，我们在 $U$ 上局部地工作，并选取局部坐标
$\{x^1,\dots,x^n\}$。设
\[
  \Omega=\ee^{\rho(x)}\dd x^1\wedge\cdots\wedge\dd x^n.
\]
局部地，
\[
  \mathrm{PV}(U)=C^{\infty}(U)[\partial_1,\dots,\partial_n],
\]
其中 $\partial_i$ 之间反交换：$\partial_i\wedge\partial_j=-\partial_j\wedge\partial_i$。

% [Tutor Note: 读上方的交换图时请注意两排箭头方向\emph{相反}：上排 de Rham
% 微分的次数是 $+1$，下排 $\Delta$ 的次数是 $-1$（把 $k$ 形式的多向量场送到
% $k-1$ 形式）。这是“多向量场是形式的对偶”这一事实的体现；若把 $\Delta$
% 的次数误记为 $+1$，BV 括号的次数（也是 $+1$）就会算错。]

% >>> p051

为方便起见，引入 Grassmann 变量 $\theta_i$ 来表示
\[
  \theta_i\equiv\partial_j,\qquad \theta_i\theta_j=-\theta_j\theta_i.
\]
于是 $\mu\in\mathrm{PV}(U)$ 可以写成
\[
  \mu=\mu(x^i,\theta_i)=\sum_I\mu_I(x^i)\theta^I.
\]
设 $\frac{\partial}{\partial\theta_i}$ 为关于 $\theta_i$ 的导数，其定义为
\[
  \frac{\partial g}{\partial\theta_i}=0,\qquad
  \frac{\partial(\theta_i g)}{\partial\theta_i}=g,
\]
这里 $g$ 不含 $\theta_i$。这些算子也是反交换的：
\[
  \frac{\partial}{\partial\theta_i}\frac{\partial}{\partial\theta_j}
  =-\frac{\partial}{\partial\theta_j}\frac{\partial}{\partial\theta_i}.
\]
例如
\[
  \begin{aligned}
    \frac{\partial}{\partial\theta_1}\frac{\partial}{\partial\theta_2}
      (\theta_1\theta_2)
    &=-\frac{\partial}{\partial\theta_1}\frac{\partial}{\partial\theta_2}
      (\theta_2\theta_1)
     =-\frac{\partial\theta_1}{\partial\theta_1}=-1,\\
    \frac{\partial}{\partial\theta_2}\frac{\partial}{\partial\theta_1}
      (\theta_1\theta_2)&=1.
  \end{aligned}
\]

% [Tutor Note: 上式第一行左端 $\theta_i\equiv\partial_j$ 的指标不匹配：按定义
% $\theta_i$ 应对偶于 $\partial_i$，即应为 $\theta_i\equiv\partial_i$。原文
% PDF p.51 作 $\theta_i\equiv\partial_j$，疑为笔误，此处照录并提示核对原书
% p.50 / PDF p.51。这个指标错位不影响后文（后文只用 $\theta_i$ 作为独立的
% 奇坐标）。]

\begin{wprop}{6.2}
在上述假设与记号下，有
\[
  \Delta=\sum_i\frac{\partial}{\partial x^i}
  \frac{\partial}{\partial\theta_i}
  +\sum_i\partial_i\rho\,\frac{\partial}{\partial\theta_i}.
\]
\end{wprop}

\begin{proof}
习题。
\end{proof}

特别地，对向量场
\[
  X=\sum_i V^i(x)\frac{\partial}{\partial x^i}
  =\sum_i V^i(x)\theta_i,
\]
我们有
\[
  \operatorname{Div}(X)=\Delta X
  =\sum_i\left(\partial_iV^i+\partial_i\rho\,V^i\right).
\]
这恰是散度的公式。

注意第一项 $\sum_i\frac{\partial}{\partial x^i}\frac{\partial}{\partial\theta_i}$
看起来像一个 Laplacian，但它\emph{不是}——$\theta_i$ 是奇变量。尽管如此，
$\Delta$ 有时被称为 \emph{BV Laplacian}，或\emph{奇 Laplacian}。

回顾积分在同调层面的观点：
\begin{figure}[htbp]
  \centering
  \includegraphics[width=0.75\linewidth]{fig_p051_1.png}
  \caption*{积分的同调观点（原文 PDF p.51 示意图，原文无编号）}
\end{figure}
由该图我们得到
\[
  \int_{dR}\sim H^n_{dR},\qquad \int_{BV}\sim H^0_{\Delta}.
\]
在量子场论中，我们必须处理 $\infty$ 维几何，即 $n=\infty$。此时 $H^n_{dR}$
失去意义，而 $H^0_{\Delta}$ 变得更好，问题便归结为在 QFT 中构造 $\Delta$。
我们将在微扰重整化理论中看到如何做到这一点。

% [Tutor Note: “奇 Laplacian”之名强调 $\Delta$ 是二阶算子且 $\Delta^2=0$，
% 但它的“二阶性”来自奇变量的一阶导数（$\partial_x\partial_\theta$），与通常
% 椭圆 Laplacian 的 ${}^2=0$ 性质完全不同。术语上不要与 Hodge Laplacian
% $\dd\dd^*+\dd^*\dd$ 混淆：后者是偶数阶、不平方为零。]


% >>> p052

\subsection{Batalin--Vilkovisky（BV）代数}

\begin{wdefn}{6.3}
\emph{BV 代数}是一对 $(A,\Delta)$，其中
\begin{itemize}[leftmargin=1.4em,itemsep=0.1em]
  \item[\textcircled{1}] $A$ 是分次交换结合代数：
        $a\cdot b=(-1)^{|a||b|}b\cdot a$，$\forall a,b\in A$；
  \item[\textcircled{2}] $\Delta:A\to A$ 是一个 $\deg=1$ 的“二阶”算子，且
        满足 $\Delta^2=0$。
\end{itemize}
\end{wdefn}

这里“二阶”的意思如下。考虑 $\Delta$ 不是导子的失效程度，用所谓的
\emph{BV 括号}来衡量：
\[
  \{a,b\}:=\Delta(ab)-(\Delta a)b-(-1)^{|a|}a\Delta b.
\]
于是 $\{-,-\}$ 对 $a$ 与 $b$ 两个分量都满足 Leibniz 法则。因此
\[
  \{-,-\}:\ A\times A\longrightarrow A
\]
是 $\deg=1$ 的，并满足

\textcircled{1}\ 分次对称性
\[
  \{a,b\}=(-1)^{|a||b|}\{b,a\}
\]

\textcircled{2}\ 分次 Leibniz 法则
\[
  \{a,bc\}=(-1)^{(1+|a|)|b|}b\{a,c\}+\{a,b\}c
\]

\textcircled{3}\ 由构造，$\{-,-\}$ 与 $\Delta$ 也相容：
\[
  \Delta\{a,b\}=-\{\Delta a,b\}-(-1)^{|a|}\{a,\Delta b\}.
\]
如果我们还有一个
\[
  Q:\ A\longrightarrow A,\qquad \deg=1,
\]
它是导子、平方为零 $Q^2=0$，且与 $\Delta$ 相容
\[
  Q\Delta+\Delta Q=0,
\]
则三元组 $\{A,Q,\Delta\}$ 称为 \emph{DGBV 代数}。

% [Tutor Note: BV 括号定义中的 $(-1)^{|a|}$ 极易放错位置：
% $\{a,b\}=\Delta(ab)-(\Delta a)b-(-1)^{|a|}a\Delta b$。若把它写成
% $(-1)^{|b|}$，分次对称性 $\{a,b\}=(-1)^{|a||b|}\{b,a\}$ 立刻不成立。
% 记忆口诀：第二个“修正项”里的符号要移到 $\Delta$ 越过 $a$ 的位置，故与 $a$
% 的次数有关。]


\begin{wexam}{6.4}
设 $\Omega$ 是 $M$ 上的体积形式，$\Delta$ 是 $\mathrm{PV}(M)$ 上相应的散度
算子。则
\[
  (\mathrm{PV}(M),\Delta)
\]
构成一个 BV 代数。相应的 BV 括号
\[
  \{-,-\}:\ \mathrm{PV}^k(M)\times\mathrm{PV}^l(M)
  \longrightarrow\mathrm{PV}^{k+l-1}(M)
\]
不依赖于 $\Omega$ 的选取，且 $\{-,-\}$ 就是 Schouten--Nijenhuis 括号（至多
差一个符号）（事实上
$\{\alpha,\beta\}=-(-1)^{|\alpha|}\{\alpha,\beta\}_{SN}$）。

在局部坐标下，设
\[
  \Omega=\ee^{\rho(x)}\dd x^1\cdots\dd x^n,
\]
则
\[
  \Delta=\sum_i\frac{\partial}{\partial x^i}
  \frac{\partial}{\partial\theta_i}
  +\sum_i\partial_i\rho\,\frac{\partial}{\partial\theta_i}.
\]
我们得到
\[
  \{x^i,\theta_j\}=\delta^i_{\ j},
  \qquad
  \{x^i,x^j\}=\{\theta_i,\theta_j\}=0.
\]
这是 Poisson 括号的奇类比。
\end{wexam}

% [Tutor Note: Schouten–Nijenhuis 括号在不同文献中的定义相差
% $(-1)^{|\alpha|}$ 或一个整体负号很常见。原文写作
% $\{\alpha,\beta\}=-(-1)^{|\alpha|}\{\alpha,\beta\}_{SN}$，即 BV 括号与 SN 括号
% 差一个与 $|\alpha|$ 有关的符号。按术语表，“Schouten bracket”译作
% “Schouten 括号”，全称“Schouten–Nijenhuis 括号”；原文排印作
% “SchoutenNijenhuis”（缺连字符），属排版笔误。核对符号时请以原书
% p.51 / PDF p.52 为准。]


% >>> p053

\begin{wrem}{6.5}
本例说明了一个关于量子化的重要观点：量子化归结为构造无穷维的体积 / 测度。
\begin{figure}[htbp]
  \centering
  \includegraphics[width=0.50\linewidth]{fig_p053_1.png}
  \caption*{量子化：从经典 Poisson 括号到 BV 算子（原文 PDF p.53 示意图，原文无编号）}
\end{figure}
\end{wrem}

% [Tutor Note: 该图是本章的概念枢纽：把经典一侧的 Poisson 括号 $\{-,-\}$ 换成
% 量子一侧的 BV 算子 $\Delta$（由体积形式 $\Omega$ 构造）。换言之，所谓
% “量子化”就是构造 $\infty$ 维的体积 / 测度，而 $\Delta$ 正是该测度的散度
% 算子。这也是第 7 章有效 BV 量子化与第 9 章 Fedosov 量子化的共同出发点。]

\begin{wexam}{6.6}
设 $X$ 是复流形，$\Omega$ 是全纯体积形式。这样的对 $(X,\Omega)$ 称为
\emph{Calabi--Yau 流形}。我们可以用下式定义多向量场的复类比
\[
  \mathrm{PV}^{k,l}(X)=\Omega^{0,l}\left(X,\wedge^k T^{1,0}_X\right).
\]
设 $\Delta$ 是与 $\Omega$ 相联系的散度算子
\[
  \Delta:\ \mathrm{PV}^{k,l}\longrightarrow\mathrm{PV}^{k-1,l}.
\]
设 $\bar{\partial}$ 是 Dolbeault 微分
\[
  \bar{\partial}:\ \mathrm{PV}^{k,l}\longrightarrow\mathrm{PV}^{k,l+1}.
\]
则 $(\mathrm{PV}^{-\bullet,\bullet},\bar{\partial},\Delta)$ 是一个 DGBV。这里
次数约定是 $\mathrm{PV}^{k,l}$ 的次数为 $l-k$。该 DGBV 控制 $X$ 的形变理论。
\end{wexam}

\begin{wexam}{6.7}
设 $G$ 是 Lie 群，Lie 代数为 $\Lie{g}$。令
\[
  A=\Sym(\Lie{g}^{\vee}[-1])\otimes\Sym(\Lie{g}[2])
  =\C[e_\alpha,c^\alpha].
\]
这里
\begin{itemize}[leftmargin=1.4em,itemsep=0.15em]
  \item $\{e_\alpha\}$ 是 $\Lie{g}$ 的一组基，$\deg(e_\alpha)=-2$。它们彼此
        交换：$e_\alpha e_\beta=e_\beta e_\alpha$。
  \item $\{c^\alpha\}$ 是 $\Lie{g}$ 的对偶基，$\deg(c^\alpha)=1$。它们彼此
        反交换：$c^\alpha c^\beta=-c^\beta c^\alpha$。
  \item $\Delta=\sum_\alpha\frac{\partial}{\partial c^\alpha}
        \frac{\partial}{\partial e_\alpha}$，且
        $\{c^\alpha,e_\beta\}=\delta^\alpha_{\ \beta}$。
\end{itemize}
则 $(A,\Delta)$ 定义一个 BV 代数。

与上面的多向量场构造相比较
\[
  \begin{array}{ccc}
    \deg=1 & c^{\alpha}\longleftrightarrow x^{i} & \deg=0\\[2pt]
    \deg=-2 & e_{\alpha}\longleftrightarrow \theta_{i} & \deg=-1
  \end{array}
\]
我们可以把 $A$ 看作奇空间 $\Lie{g}[1]$ 上的多向量场
\[
  A=\mathrm{PV}(\Lie{g}[1])=\mathcal{O}\left(\Lie{g}[1]\oplus
  \Lie{g}^{\vee}[-2]\right).
\]
\end{wexam}

\begin{wexam}{6.8}
一般地，设 $W$ 是分次向量空间。记
\[
  T^*W[-1]=W\oplus W^{\vee}[-1].
\]
则
\[
  \mathrm{PV}(W):=\mathcal{O}(T^*W[-1])
  =\Sym(W^{\vee})\otimes\Sym(W[1]),
\]
带有典范的 $\deg=1$ 的 BV 括号 $\{-,-\}$。在特殊情形
\begin{itemize}[leftmargin=1.4em,itemsep=0.1em]
  \item $W=\R^n$ $\Longrightarrow$ $\mathrm{PV}(\R^n)$；
  \item $W=\Lie{g}[1]$ $\Longrightarrow$ $\mathrm{PV}(\Lie{g}[1])$。
\end{itemize}
特别地，为了理解 $\int_{V\times\Lie{g}[1]}$，我们将需要考察
\[
  \mathrm{PV}(V\times\Lie{g}[1]).
\]
\end{wexam}

% >>> p054

\subsection{经典 BV 主方程（Classical BV Master Equation）}

现在我们讨论所谓的 \emph{BV 主方程}，它是经典与量子规范理论的基本关系。

\begin{wdefn}{6.9}（经典 BV 主方程）
设 $(A,\Delta)$ 是 BV 代数，$\{-,-\}$ 是 BV 括号。一个次数为 $0$ 的元素
$S_0\in A$ 称为满足\emph{经典 BV 主方程}，若
\[
  \{S_0,S_0\}=0\qquad\text{(CME)}.
\]
\end{wdefn}

$S_0$ 通过
\[
  \delta=\{S_0,-\}:\ A\longrightarrow A
\]
在 $A$ 上诱导一个算子。由于 BV 括号的次数为 $1$、$S_0$ 的次数为 $0$，算子
$\delta$ 的次数为 $1$。并且由 $\{-,-\}$ 的 Leibniz 法则，$\delta$ 是一个导子。

\begin{wprop}{6.10}
CME $\{S_0,S_0\}=0$ 蕴含
\[
  \delta^2=0.
\]
\end{wprop}

\begin{proof}
由 Jacobi 恒等式
\[
  \{\{S_0,S_0\},-\}
  =-2\{S_0,\{S_0,-\}\}
  =-2\delta^2.
\]
\end{proof}


\begin{wexam}{6.11}（对称性）
设 $G$ 是 Lie 群，Lie 代数为 $\Lie{g}$。记 $\{e_\alpha\}$ 为 $\Lie{g}$ 的一组
基，$\{c^\alpha\}$ 为 $\Lie{g}^{\vee}$ 的对偶基。设 $(A,\Delta)$ 是
$\Lie{g}[1]$ 上多向量场的 BV 代数
\[
  A=\Sym(\Lie{g}^{\vee}[-1])\otimes\Sym(\Lie{g}[2])
  =\C[e_\alpha,c^\alpha],\qquad
  \deg(c^\alpha)=1,\quad \deg(e_\alpha)=-2,
\]
且
\[
  \Delta=\sum_\alpha\frac{\partial}{\partial c^\alpha}
  \frac{\partial}{\partial e_\alpha},
  \qquad
  \{c^\alpha,e_\beta\}=\delta^\alpha_{\ \beta}.
\]
在所选基下 Lie 括号为
\[
  [e_\alpha,e_\beta]=\sum_\gamma f^{\gamma}_{\ \alpha\beta}e_\gamma,
\]
其中 $f^{\gamma}_{\ \alpha\beta}$ 是 $\Lie{g}$ 的结构常数。则次数为 $0$ 的
元素
\[
  I=\frac{1}{2}\sum_{\alpha,\beta,\gamma}
  f^{\gamma}_{\ \alpha\beta}c^\alpha c^\beta e_\gamma\in A
\]
满足经典主方程
\[
  \{I,I\}=0.
\]
让我们考虑诱导出的微分
\[
  \delta=\{I,-\}:\ A\longrightarrow A.
\]
\end{wexam}

% [Tutor Note: 本页原文两处笔误：(i) $\deg(e_\alpha=-2)$ 括号错位，应为
% $\deg(e_\alpha)=-2$；(ii) 正文另有一处写作 $\{c^\alpha,e_\alpha\}
% =\delta^\alpha_{\ \beta}$，左端两个指标都是 $\alpha$ 而右端是 $\alpha,\beta$，
% 应为 $\{c^\alpha,e_\beta\}=\delta^\alpha_{\ \beta}$。两处均照录原意并在
% 译文中按正确指标排版，请核对原书 p.53 / PDF p.54。]

% >>> p055

由于 $A=\C[c^\alpha,e_\alpha]$ 是自由代数，我们只需知道 $\delta$ 如何作用在
生成元上。不难看出
\[
  \delta(c^\gamma)=\frac{1}{2}f^{\gamma}_{\ \alpha\beta}c^\alpha c^\beta,
  \qquad
  \delta(e_\alpha)=\frac{1}{2}f^{\gamma}_{\ \alpha\beta}c^\beta e_\gamma.
\]
因此导子 $\delta$ 由向量场
\[
  \delta=\frac{1}{2}\sum_{\alpha,\beta,\gamma}
  f^{\gamma}_{\ \alpha,\beta}c^\alpha c^\beta
  \frac{\partial}{\partial c^\gamma}
  +\sum_{\alpha,\beta,\gamma}
  f^{\gamma}_{\ \alpha\beta}c^\beta e_\gamma
  \frac{\partial}{\partial e_\alpha}
\]
给出。若把 $\delta$ 限制到 $\C[c^\alpha]=\Sym(\Lie{g}^{\vee}[-1])$，则
$\delta=d_{\mathrm{CE}}$。因此
\[
  \left(\Sym(\Lie{g}^{\vee}[-1]),\delta\right)=C^{\bullet}(\Lie{g}).
\]
在完整的 BV 代数 $A$ 上，我们有
\[
  (A,\delta)=C^{\bullet}(\Lie{g},\Sym(\Lie{g}[2])),
\]
其中 $\Lie{g}$ 通过伴随表示作用在 $\Sym(\Lie{g}[2])$ 上。

% [Tutor Note: 本页三处 $\frac{1}{2}$ 的记账需要核对：$\delta(c^\gamma)
% =\frac{1}{2}f^{\gamma}_{\ \alpha\beta}c^\alpha c^\beta$ 与紧随其后的向量场
% 第一项一致；但 $\delta(e_\alpha)$ 被写成
% $\frac{1}{2}f^{\gamma}_{\ \alpha\beta}c^\beta e_\gamma$，而向量场第二项
% $\sum f^{\gamma}_{\ \alpha\beta}c^\beta e_\gamma\partial/\partial e_\alpha$
% 却没有 $\frac{1}{2}$。按 BV 括号直接计算（也可由下方向量场反推），应为
% $\delta(e_\alpha)=f^{\gamma}_{\ \alpha\beta}c^\beta e_\gamma$（无 $\frac{1}{2}$），
% 故第一个式子的 $\frac{1}{2}$ 疑为原文笔误。按体例照录原文，请核对原书
% p.54 / PDF p.55。]

\begin{wexam}{6.12}（对称性 $+$ 物质）
设 $G$ 是作用于流形 $X$ 的 Lie 群，$X$ 的坐标为 $\{x^i\}$，$\{e_\alpha\}$ 是
Lie 代数 $\Lie{g}$ 的一组基。无穷小作用给出一个 Lie 代数同态
\[
  \begin{aligned}
    \Lie{g}&\longrightarrow\operatorname{Vect}(X),\\
    e_\alpha&\longmapsto V_\alpha
    =V^i_\alpha(x)\frac{\partial}{\partial x^i}.
  \end{aligned}
\]
考虑 $G$-不变函数
\[
  f:\ X\longrightarrow\C,
\]
它模拟一个规范不变的经典作用量。考虑
\[
  A=\mathrm{PV}(X\times\Lie{g}[1])=\mathrm{PV}(X)\otimes\mathrm{PV}(\Lie{g}[1])
  =\C[x^i,\theta_i,e_\alpha,c^\alpha],
\]
其中 $\theta_i=\partial_i$，$\{c^\alpha\}$ 是对偶基。我们在下表中总结这些变量。

\begin{center}
\normalfont\small
\begin{tabular}{lcccc}
  \toprule
  变量 & $c^\alpha$ & $x^i$ & $\theta_i$ & $e_\alpha$ \\
  \midrule
  次数（Degree） & $1$ & $0$ & $-1$ & $-2$ \\
  物理意义（Physics） & 鬼场 & （物质）场 & 反场 & 反鬼场 \\
  \bottomrule
\end{tabular}
\end{center}

我们可以把函数 $I=\frac{1}{2}f^{\gamma}_{\ \alpha\beta}c^\alpha c^\beta e_\gamma$
通过纳入 $x$ 与 $G$-不变函数 $f$ 延拓为
\[
  S_0=f(x)+c^\alpha V^i_\alpha(x)\theta_i
  +\frac{1}{2}f^{\gamma}_{\ \alpha\beta}c^\alpha c^\beta e_\gamma.
\]
\end{wexam}

\begin{wprop}{6.13}
$S_0$ 满足经典主方程。
\end{wprop}

\begin{proof}
习题。
\end{proof}

% [Tutor Note: 上表按原文重排（booktabs）。原文表头第一格为空，第一列是
% “Degree / Physics”两行；本译本把变量名提到表头、把“次数”“物理意义”放到
% 第一列，内容与原文一一对应。术语依据：ghost 鬼场、anti-field 反场、
% anti-ghost 反鬼场（术语表）。]

% [Tutor Note: Prop.\ 6.13 的“Proof: Exercise.”省略了实质验证：把
% $S_0=f+c^\alpha V^i_\alpha\theta_i+\frac{1}{2}f^{\gamma}_{\ \alpha\beta}
% c^\alpha c^\beta e_\gamma$ 代入 $\{S_0,S_0\}=0$，需要同时用到
% (i) $f$ 的 $G$-不变性 $V_\alpha(f)=0$；(ii) 结构常数的 Jacobi 恒等式；
% (iii) $\Lie{g}\to\operatorname{Vect}(X)$ 是 Lie 代数同态，即
% $[V_\alpha,V_\beta]=f^{\gamma}_{\ \alpha\beta}V_\gamma$。三者缺一不可，
% 而原文全部省略。]
